
This work validates the experimental protocol for measuring calibration quality in code generation tasks. The simulation demonstrates that temperature scaling optimization converges reliably, that ECE can be computed stably with 15-bin partitioning, and that ranking-based accuracy is preserved as expected from theoretical properties of monotonic transformations.

The key contribution is not empirical findings about code generation models, but rather a validated and operational experimental pipeline. Future work can use this protocol to:
\begin{enumerate}
\item Measure baseline calibration quality of Code Llama 7B (and other models)
\item Quantify the effectiveness of temperature scaling on real model outputs
\item Compare calibration methods (temperature vs. vector vs. matrix scaling)
\item Establish whether code generation exhibits different calibration characteristics than classification
\end{enumerate}

The simulation approach enables rapid validation of experimental infrastructure before committing computational resources to multi-hour production experiments. It also clearly delimits technical validation (what this work accomplishes) from empirical measurement (what future work is needed).

Calibration research for code generation remains in its earliest stages. Before investigating advanced questions about confidence-based resource allocation or adaptive iteration control, the field needs basic empirical measurements: What is the baseline ECE for representative code generation models? Does temperature scaling improve calibration on real code generation distributions? This work provides the validated experimental protocol needed to answer these questions.
